% IMMUTABLE RESEARCH REPORT::FINAL-LINK-LIT::2026-10-08_003222
\documentclass[11pt]{article}
\usepackage[T1]{fontenc}
\usepackage[margin=0.88in]{geometry}
\usepackage{amsmath,amssymb,amsthm,mathtools,booktabs,array,longtable,hyperref,xcolor}
\setlength{\emergencystretch}{3em}
\hypersetup{colorlinks=true,linkcolor=blue,citecolor=blue,urlcolor=blue}

\newtheorem{theorem}{Theorem}
\newtheorem{lemma}[theorem]{Lemma}
\newtheorem{proposition}[theorem]{Proposition}
\newtheorem{definition}[theorem]{Definition}
\newcommand{\SAT}{\operatorname{SAT}}
\newcommand{\poly}{\operatorname{poly}}
\newcommand{\COVER}{\textsc{Cover-Shrink}}
\newcommand{\Known}{\textsc{Known Theorem}}
\newcommand{\Open}{\textsc{Open}}
\newcommand{\Conditional}{\textsc{Conditional}}
\newcommand{\Restricted}{\textsc{Restricted Positive}}

\title{FINAL-LINK-LIT:\\The Exact Literature Frontier for a Charged SAT Cover-Shrinker}
\author{RESEARCH agent report}
\date{8 October 2026}

\begin{document}
\maketitle

\begin{abstract}
This report tests the final proposed WRAPPER link against primary literature
on sparsification, Turing compression, kernelization, backdoors,
representative families, split-and-list, universal sets, covering codes,
knowledge compilation, semantic hashing, selectors, branching, and
Circuit-SAT self-improvement.  The clean target is an exact all-input
constructor that, on an ordinary CNF of encoded length $M$, emits at most
$M^{1/3+o(1)}$ ordinary SAT instances of length at most
$M^{2/3+o(1)}$, with total discovery, materialization, replay, and local work
$M^{1+o(1)}$.  The most important finding is quantitative: on dense
$n$-variable $3$-CNF, this target is the equality case $b+c=3$ of the
Dell--van Melkebeek multi-query sparsification frontier.  Their theorem rules
out $b+c<3$ unless $\mathrm{coNP}\subseteq\mathrm{NP/poly}$, but neither
rules out nor constructs the balanced point $(b,c)=(1,2)$.  Backdoor DNFs
give the closest exact syntax for the proposed forward cover, while known
detection bounds and residual-size guarantees do not meet the target on all
CNFs.  The strongest source-checkable candidate is therefore a
\emph{size-charged backdoor-DNF cover} with a structural coverage
certificate and output-sensitive materialization.  Its universal existence
and cheap acquisition remain \textbf{OPEN}.  No reviewed result completes a
proof that $P\ne NP$.
\end{abstract}

\section{Question, accounting convention, and answer}

Let $F$ be an explicitly encoded CNF of bit length $M$.  The desired forward
object must establish an exact identity of the form
\begin{equation}
  \SAT(F)=\Phi\bigl(\SAT(F_1),\ldots,\SAT(F_q)\bigr),
  \tag{1}
\end{equation}
where $\Phi$ is a source-checkable combiner, each $F_i$ is an ordinary
explicit SAT instance, and
\begin{align}
 q&\le M^{1/3+o(1)},
 & |F_i|&\le M^{2/3+o(1)}, \tag{2}\\
 T_{\rm discover}+T_{\rm write}+T_{\rm replay}+T_{\rm combine}
   &\le M^{1+o(1)}.
 && \tag{3}
\end{align}
The target permits OR, AND, or a declared bounded-depth combination, but it
does not permit a hidden pointer to the entire parent inside a nominally
short ordinary query.  If a child is a succinct source-relative object, the
oracle problem and its cost model have changed and a separate return theorem
is required.

The answer from the literature is precise.  Equation (2) is not contradicted
by the strongest known multi-query compression lower bound at the dense
$3$-SAT scale.  It lies exactly at that lower bound's unexcluded equality
frontier.  No published construction reaches (2)--(3) on every CNF.  Several
published mechanisms establish one or two of exactness, small branching,
small residuals, or cheap discovery under additional structural promises,
but none establishes all four universally.

\section{The multi-query frontier is exactly the requested dense point}

\begin{theorem}[Dell--van Melkebeek, Corollary 1]\label{thm:dvm}
Fix an integer $d\ge3$ and nonnegative reals $b,c$ with $b+c<d$.  Unless
$\mathrm{coNP}\subseteq\mathrm{NP/poly}$, there is no polynomial-time
reduction from $n$-variable $d$-SAT to an arbitrary language that makes at
most $O(n^b)$ queries, each of bit length $O(n^c)$
\cite{dellvanmelkebeek2014}.
\end{theorem}

The theorem charges the number of queries and the maximum query length.  It
is therefore more directly relevant than a one-output kernel lower bound.
For a dense $3$-CNF, suppressing logarithmic encoding factors gives
$M=n^{3+o(1)}$.  The desired bounds become
\[
 q=M^{1/3+o(1)}=n^{1+o(1)},\qquad
 |F_i|=M^{2/3+o(1)}=n^{2+o(1)}.
\]
Thus $(b,c)=(1,2)$ and
\[
 b+c=3=d. \tag{4}
\]
The target is not in the conditionally forbidden region $b+c<d$.  It is on
the equality boundary.  This observation matters in both directions.  The
known lower bound does not dispose of the proposal, but equality is not an
upper-bound theorem.  The trivial one-query reduction with $(b,c)=(0,d)$
also lies on this boundary.  Dell--van Melkebeek do not provide a balanced
$n$-by-$n^{d-1}$ construction.

In aggregate-output language, the target satisfies
\[
 q\max_i|F_i|=M^{1+o(1)}. \tag{5}
\]
It conserves, rather than compresses below, the parent-scale number of bits.
This explains why generic OR-compression lower bounds do not automatically
exclude it.  The missing content is a semantic redistribution of the source
into individually smaller exact questions without paying for a successful
route in advance.

\subsection{The sparse regime is strictly more demanding}

For constant-width CNF with $m$ clauses, an explicit bit encoding has
$M=\widetilde\Theta(n+m)$.  When $m=\Theta(n)$, (2) requests children of
length $n^{2/3+o(1)}$.  Such a child cannot even name all $n$ source
variables explicitly.  A valid constructor must eliminate variables,
factor the instance, or derive a decision-preserving quotient.  Clause
sparsification alone is insufficient.  Source pointers can avoid rewriting,
but a pointer is not an ordinary short CNF query.  Consequently a theorem
proved only for dense $M=n^{3+o(1)}$ inputs would not establish the stated
all-$M$ WRAPPER contract.

\section{Quantitative sparsification does not interpolate to the target}

The Impagliazzo--Paturi--Zane sparsification lemma is constructive and exact.
In the quantitative Calabro--Impagliazzo--Paturi formulation, for fixed
$d$ and $\varepsilon>0$, an $n$-variable $d$-CNF is the OR of at most
$2^{\varepsilon n}$ $d$-CNFs, each with
\begin{equation}
 O\!\left(\left(\frac d\varepsilon\right)^{3d}n\right)
 \tag{6}
\end{equation}
clauses; enumeration has $n^{O(d)}$ delay
\cite{ipz2001,calabroimpagliazzopaturi2010}.  The 2026 measure-and-conquer
proof gives the same structural lemma with the best known dependence
$C_{\varepsilon,d}=(d/\varepsilon)^{O(d)}$ on variable occurrence
\cite{lokshtanov2026}.

For $d=3$, forcing polynomial fanout $q\le n^{1+o(1)}$ requires
\[
 \varepsilon\le\frac{(1+o(1))\log_2 n}{n}.
\]
Substitution in (6) gives
\[
 |F_i|_{\rm clauses}
 =O\!\left(n\left(\frac{3n}{\log n}\right)^9\right)
 =O\!\left(\frac{n^{10}}{\log^9 n}\right). \tag{7}
\]
This is far above $n^2$.  Conversely, making (6) at most $n^2$ requires
$\varepsilon=\Omega(n^{-1/9})$, which gives
$q=2^{\Omega(n^{8/9})}$.  The newer
$(d/\varepsilon)^{O(d)}$ bound improves constants and proof structure, but
when $\varepsilon=\Theta(\log n/n)$ it still yields a child size
$n^{1+\Omega(d)}$ rather than the requested $n^2$ guarantee.

This calculation identifies the exact failure mode.  Fixed-$\varepsilon$
sparsification gives small children and exponentially many of them.
Vanishing $\varepsilon$ gives polynomially many children but loses the child
bound.  The final link requires a new balanced lemma, not a different
constant in the standard sparsification proof.

Santhanam and Srinivasan prove genuine non-sparsifiability theorems for
broader CNFs.  In particular, for every superlinear $f(n)$ there are CNFs of
size at most $f(n)$ that are not ORs of $2^{n-\varepsilon n}$ linear-size
CNFs, and for constants $k$ and $k'>2k$ they obtain a hierarchy separating
$n^{k'}$-size CNFs from subexponential ORs of $n^k$-size CNFs
\cite{santhanamsrinivasan2011}.  Their paper explicitly notes that the first
argument does not settle whether $n^3$-size CNFs sparsify to $n^2$-size
CNFs, and the hierarchy condition $k'>2k$ does not cover $(k',k)=(3,2)$.
Hence this line also leaves the balanced dense target open.

\section{The closest exact syntax: a size-charged backdoor DNF}

Ordyniak, Schidler, and Szeider define a $\mathcal C$-backdoor DNF as a
tautological DNF of terms corresponding to partial assignments
$\tau_1,\ldots,\tau_q$ such that every restricted formula $F[\tau_i]$
belongs to a tractable base class $\mathcal C$
\cite{ordyniakschidlerszeider2021}.  The tautology condition gives the exact
identity
\begin{equation}
 \SAT(F)=\bigvee_{i=1}^{q}\SAT(F[\tau_i]). \tag{8}
\end{equation}
Unlike a strong backdoor set, a backdoor DNF need not enumerate every
assignment to one fixed variable set.  It can therefore be strictly smaller
than both a backdoor set and a backdoor tree.  This is the closest published
formalism to the proposed forward-only inner WRAPPER.

The paper proves fixed-parameter tractability of detecting a $k$-term
backdoor DNF for Horn, AntiHorn, $2$CNF, $2$CNF$\cup$Horn, and
$2$CNF$\cup$AntiHorn.  Its displayed construction through a supplied
variable set $B$ includes the factor
\begin{equation}
 2^{3|B|+1}+3^{|B|}|F|^{O(1)}, \qquad |B|\le k-1, \tag{9}
\end{equation}
together with a branching factor bounded as a function of $k$.  FPT in $k$
does not mean near-linear when $k=M^{1/3}$.  The theorem also does not say
that every formula has such a cover with $k\le M^{1/3}$, or that each
residual has length at most $M^{2/3}$.

Classical strong backdoors expose the same accounting issue.  Nishimura,
Ragde, and Szeider give $O(2^kN)$ strong-Horn-backdoor detection and
$O(3^kN)$ strong-$2$CNF-backdoor detection for input scale $N$
\cite{nishimuraragdeszeider2004}.  If $2^k=M^{1/3}$, these become
$M^{4/3+o(1)}$ and
\[
 M^{1+\log_2(3)/3+o(1)}\approx M^{1.5283}, \tag{10}
\]
respectively, before accounting for residual writing.  Given the backdoor,
enumeration produces $2^k$ exact tractable residuals, but a restriction can
leave $\Theta(M)$ symbols.  Small branching therefore does not imply power
shrinkage.

Recursive backdoors formalize the proposed recursive wrapper more closely:
after assignments, disconnected residual components are solved recursively.
For the empty base class, published detection and evaluation have
$2^{2^{O(k)}}|F|$ dependence, and no all-input bound on $k$ is supplied
\cite{mahlmannsiebertzvigny2021}.  The construction is exact, but its known
parameter dependence cannot be inserted at a growing polynomial value of
$k$.

\begin{definition}[Size-charged backdoor-DNF cover]\label{def:scbdnf}
A size-charged cover for an explicit CNF $F$ of length $M$ consists of
partial assignments $\tau_1,\ldots,\tau_q$ and a structural coverage
certificate such that
\begin{align*}
 \bigvee_{i=1}^q \tau_i&\equiv 1,
 &q&\le M^{1/3+o(1)},\\
 |F[\tau_i]|&\le M^{2/3+o(1)},
 &\sum_i |F[\tau_i]|&\le M^{1+o(1)},
\end{align*}
and the cover, residuals, and certificate are discovered and replayed in
$M^{1+o(1)}$ total work and permitted workspace.
\end{definition}

Definition~\ref{def:scbdnf} is a faithful literature-backed statement of
the missing final link.  If the residuals are in a tractable class, it is
already a solver.  If they are arbitrary smaller CNFs, it is the desired
forward cover for the outer wrapper.  The coverage certificate should be a
decision tree, disjoint cube partition, or another locally checkable object;
checking whether an arbitrary DNF is tautological is itself coNP-complete.
No reviewed paper proves that every CNF has such a cover.

\section{What a double wrapper can and cannot certify}

A natural double-wrapper proposal has an inner nonrecursive phase emit
candidate terms $\tau_i$ and an outer phase later obtain SAT answers for the
residuals.  This successfully separates route construction from route
selection.  Exactness follows from (8) once coverage is certified.  It avoids
asking which route is successful during construction.

One may try to verify coverage by asking whether the complement of the DNF
is satisfiable.  For terms $\tau_i$, noncoverage is
\[
 \neg\bigvee_i\tau_i=\bigwedge_i\neg\tau_i,
\]
which is a CNF whose clauses are the negated terms.  This is an exact
verification query, but unless its encoding is itself at most
$M^{2/3+o(1)}$ it reintroduces the current-scale call that the forward
wrapper was meant to avoid.  If the proposed terms arise as leaves of a
complete binary decision DAG, coverage is syntactic and needs no SAT query,
but the number of leaves can be exponential.  Thus the double wrapper fixes
causal order, not coverage size or acquisition cost.

The strongest defensible version is consequently conditional:
\begin{proposition}[Conditional forward bridge]\label{prop:bridge}
If every CNF admits a uniformly constructible size-charged cover from
Definition~\ref{def:scbdnf}, then the nonrecursive inner wrapper can emit all
ordinary child CNFs before any child SAT answer is known, and the outer
wrapper can combine those answers by OR with total forward work
$M^{1+o(1)}$.
\end{proposition}
The proposition is immediate from the definition.  Its premise, not its
composition, is the open theorem.

\begin{proposition}[Why the universal cover is the breakthrough]\label{prop:recurrence}
Suppose the cover in Definition~\ref{def:scbdnf} is available unconditionally
for every CNF, with uniform control of all $o(1)$ terms.  Recursively applying
it gives
\[
 T(M)\le M^{1/3+o(1)}T\!\left(M^{2/3+o(1)}\right)+M^{1+o(1)},
 \tag{R}
\]
and hence $T(M)=M^{1+o(1)}$.
\end{proposition}
\begin{proof}
At recursion level $j$, each instance has length
$M^{(2/3)^j+o(1)}$.  The product of all preceding branching factors makes
the number of instances $M^{1-(2/3)^j+o(1)}$.  Their aggregate length and
aggregate local work are therefore $M^{1+o(1)}$.  There are
$O(\log\log M)$ nonconstant scales, which is absorbed by the subpolynomial
factor.
\end{proof}

This proposition does not refute the cover.  It explains why a universal
unconditional construction would already be a near-linear SAT algorithm.
In a contradiction argument beginning with $P=NP$, using the assumed
$M^K$ solver to discover the cover merely reintroduces $K$; the new content
must be an acquisition procedure with the charged bound in (3).

\section{Kernelization and compression: barriers and positive clues}

Fortnow and Santhanam rule out generic OR-distillation that maps an arbitrary
number of SAT instances of maximum size $n$ to one output of size
$\poly(n)$ independent of the number of inputs, unless
$\mathrm{NP}\subseteq\mathrm{coNP/poly}$
\cite{fortnowsanthanam2011}.  Drucker obtains stronger classical and quantum
instance-compression consequences under related hypotheses
\cite{drucker2015}.  Cross-composition turns such OR-compositions into
polynomial-kernel lower bounds for parameterized problems
\cite{bodlaenderjansenkratsch2014}.  These results warn against compressing a
large alternative family into one generic short instance.

They do not directly exclude Definition~\ref{def:scbdnf}.  Its aggregate
output may be $M^{1+o(1)}$, rather than independent of the source or the
number of alternatives, and it retains multiple queries.  Theorem
\ref{thm:dvm} is the sharper applicable barrier, and the desired exponents
saturate rather than violate it.  Moreover these lower bounds are
conditional.  They cannot be used as premises inside an unconditional
$P\ne NP$ proof.

Jansen and Pieterse give a useful positive clue rather than a transfer.
For each fixed $d\ge4$, $d$-NAE-SAT has a polynomial-time kernel with
$O(n^{d-1})$ clauses and $O(n^{d-1}\log n)$ bits.  It is obtained by a
linear-algebraic representation related to Lov\'asz's bound for critically
$3$-chromatic hypergraphs, and the exponent is conditionally optimal
\cite{jansenpieterse2017}.  This is a genuine source-checkable deletion rule:
linear dependence identifies a constraint that can be removed while
preserving NAE satisfiability.

The rule does not transfer to ordinary CNF-SAT.  The standard shared-control
conversion from a $d$-clause to NAE raises arity to $d+1$.  Applying the NAE
kernel then returns $O(n^d)$ constraints, matching the original dense scale
rather than shrinking it.  For $d\ge4$, even the native NAE bound
$n^{d-1}$ exceeds the particular child target $n^{2d/3}$.  Nevertheless the
mechanism is an important design clue: an algebraic invariant can make
representative deletion exact at a compression frontier.  The missing SAT
analogue would need a sound invariant for disjunctive solution existence,
not merely for clause-vector dependence.

\section{Representative families identify a required algebraic structure}

For a rank-$r=p+q$ linear matroid and an explicit family $\mathcal S$ of
$t$ independent $p$-sets, Fomin, Lokshtanov, Panolan, and Saurabh compute a
$q$-representative subfamily of size at most
\[
 \binom{p+q}{p}. \tag{11}
\]
Their Theorem 1.1 uses
\begin{equation}
 O\!\left(\binom{p+q}{p}\,t\,p^\omega
   +t\binom{p+q}{q}^{\omega-1}\right) \tag{12}
\end{equation}
field operations, with the stated matrix-multiplication exponent
$\omega<2.373$ \cite{fominrepresentative2016}.  Representation preserves the
existence of a future disjoint $q$-set that extends a current independent
set.

This is exactly the kind of future-extension theorem a WRAPPER merge needs,
but only after a problem-specific matroid and extension predicate have been
identified.  Arbitrary CNF satisfiability under future restrictions is not a
linear-matroid independence predicate.  Encoding it as one would constitute
the missing semantic compiler.  The binomial factor must itself fit the
$M^{1/3+o(1)}$ frontier, and the explicit family size $t$ in (12) is charged.
Product-family techniques avoid listing a Cartesian product when the needed
product structure is already present; they do not supply that structure for
general CNF.  Representative sets are therefore a promising implementation
tool after a low-rank invariant is proved, not a universal cover theorem.

\section{Universal sets, covering codes, and split-and-list}

Universal sets and splitters guarantee that every assignment on every
chosen $k$-coordinate subset appears in a small explicit family
\cite{naorschulmansrinivasan1995}.  Their size is polynomial in the ambient
dimension only when the relevant coordinate dimension is logarithmic.
Applied to SAT without a theorem localizing every witness to
$k=O(\log M)$ relevant coordinates, they revert to exponentially many
patterns.

Covering codes support the best-known deterministic form of Sch\"oning's
local search.  Moser and Scheder solve Promise-Ball-$k$-SAT in
$(k-1+\varepsilon)^r$ time and obtain deterministic $k$-SAT time
\[
 O^*\!\left(\left(\frac{2(k-1)}{k}+\varepsilon\right)^n\right)
 \tag{13}
\]
for every fixed $\varepsilon>0$ \cite{moserscheder2011}.  The covering code
ensures proximity to some satisfying assignment, but the number of centers
and local-search work remain exponential in the unrestricted witness
dimension.  The method does not emit a polynomial family of individually
short ordinary SAT instances.

Split-and-list divides variables into two halves and constructs
$2^{n/2}$ partial-assignment vectors, reducing SAT-like compatibility to an
orthogonal-vector or meet-in-the-middle problem
\cite{patrascuwilliams2010}.  It gives exact shared combination structure,
and modern equality-matrix decompositions improve the cost of combining
lists in some regimes.  It does not make the lists polynomial in $M$ when
$n$ is superlogarithmic.  Its useful lesson for WRAPPER design is to share
combination work across children; it does not supply the required small
frontier.

\section{Knowledge compilation and semantic hashing occur after acquisition}

Reduced ordered BDDs are canonical for a fixed variable order, so exact
hash-consing can merge equal subgraphs by the tuple consisting of a variable
and its two child identifiers \cite{bryant1986}.  Decision-DNNF and related
languages permit stronger decomposition.  Recent semantic hashing can
nominate equivalent subcircuits of an already compiled decision-DNNF and
validate candidates with provenance-aware exact checks
\cite{bierelagniezlonca2026}.  Tree Decision Diagrams similarly give
canonical and fixed-parameter compilation routes for bounded-width sources
\cite{capelli2026}.

These are valid restricted WRAPPER routes.  Amarilli, Monet, and Senellart
give a particularly explicit charged compiler: from a Boolean circuit and a
supplied width-$k$ tree decomposition $T$, they construct an equivalent
d-SDNNF in
\[
 O\!\left(|T|2^{(4+\varepsilon)k}\right)
\]
time for every fixed $\varepsilon>0$ \cite{amarillimonetsenellart2018}.
Thus logarithmic width can produce a polynomially charged compiled route,
with constants still relevant to (3).  If the source has bounded treewidth,
incidence width, or a favorable variable order, exact merging can therefore
be cheap.
On general CNFs, compilation may be exponential.  Bova, Capelli, Mengel, and
Slivovsky construct polynomial-size CNFs whose equivalent DNNFs require
strongly exponential size \cite{bovaetal2016}.  This lower bound concerns
full function representation, so it does not rule out an equisatisfiable
multi-query cover.  It does rule out the generic plan ``compile every CNF
fully, then hash and minimize'' as the near-linear acquisition phase.

Syntactic Weisfeiler--Leman hashing and exact isomorphism checks can merge
renamed duplicate residuals, while semantic hashes can accelerate candidate
selection.  Neither changes the number of semantic classes.  Hashes belong
after a coverage or compilation theorem, not in place of one.  Changing a
four-digit hash to six or eight digits changes finite collision work but not
the exponents in (2).

Valiant--Vazirani isolation gives a complementary warning.  Pairwise
independent hashes add random XOR constraints to isolate a witness with
inverse-polynomial probability \cite{valiantvazirani1986}.  The constrained
child still contains essentially all of $F$ plus the hash constraints.
Enumerating every hash bucket for deterministic exact coverage creates
$2^r$ children of length $\Theta(M)$ when the hash has $r$ output bits;
sampling fewer buckets is randomized and does not give an exact cover.
Changing $r$ therefore trades isolation probability against branch count,
but does not supply the missing encoded-size contraction.

\section{Selectors and self-reducibility expose circular route choice}

A polynomial-time selector for a language returns one of two inputs and is
guaranteed to return a member if either input is a member.  Buhrman and
Torenvliet prove that every p-selective self-reducible set lies in $P$, with
the converse giving a characterization of $P$
\cite{buhrmantorenvliet1996}.  SAT is self-reducible.  Therefore a cheap rule
that, at every Shannon branch, keeps one satisfiable child whenever one
exists is already solver-strength information.

Under the assumption $P=NP$, such a selector exists qualitatively.  This
does not control its polynomial exponent, workspace, or query size.  If it
is implemented by the same $M^K$ SAT decider whose exponent the WRAPPER is
trying to remove, the construction inherits $K$.  The forward cover avoids
this circularity only by retaining all certified alternatives and letting
the outer phase receive their answers later.  That is why Definition
\ref{def:scbdnf} asks for a small exhaustive cover, not a successful-route
selector.

\section{Circuit-SAT self-improvement is a precedent, not the missing compiler}

Williams proves that if Circuit-SAT on $n$-input circuits of size
$2^{\alpha n+o(n)}$ runs in $2^{\beta n+o(n)}$ time, for fixed positive
$\alpha\le\beta$, then Circuit-SAT on subexponential-size circuits runs in
$2^{(\beta-\alpha)n+o(n)}$ time \cite{williams2024}.  The proof packages an
inner solver into a forward circuit while leaving outer variables exposed.
This is a rigorous analogue of the double-wrapper architecture.

The theorem does not instantiate Definition~\ref{def:scbdnf} in the
ordinary-CNF polynomial-input regime.  Williams explicitly asks whether the
self-improvement phenomenon reaches the $P=NP$ setting and records
bottlenecks in the current argument.  Taking an assumed $M^K$ solver on a
large-circuit scale improves a fixed exponential rate; it does not uniformly
manufacture children of encoded length $M^{2/3}$ at linear aggregate cost
for every CNF.  The literature therefore supports the architecture but not
the final all-input compiler.

\section{A source-checkable construction program}

The literature suggests a concrete portfolio, with exact failure points.
The constructor may first apply polynomial-time normalizations and component
factorization.  It may then search, under explicit time caps, for a bounded
treewidth decomposition, a small strong backdoor, a compact backdoor DNF,
an algebraic representative rule, or a small exact decision diagram.  Every
successful member returns a syntactic certificate, its residual children,
and a combiner.  Failed attempts are charged and discarded.  Exactness of a
finite portfolio is immediate on the union of its promise domains.

What remains is not selection among successful members.  It is a universal
coverage theorem asserting that at least one member succeeds on every CNF
with the simultaneous bounds (2)--(3).  Sparse high-width formulas, formulas
with large minimum backdoors, and formulas with exponentially large compiled
representations are not covered by the cited positive theorems.

The most useful new target can be divided into three lemmas.

\paragraph{Dense frontier lemma.}
For every $n$-variable $3$-CNF of length $n^{3+o(1)}$, construct in
$n^{3+o(1)}$ work an exact OR of at most $n^{1+o(1)}$ ordinary CNFs, each of
length at most $n^{2+o(1)}$.  This is the $(b,c)=(1,2)$ equality point left
open by Theorem~\ref{thm:dvm}.

\paragraph{Sparse elimination lemma.}
For every constant-width CNF of length $M=n^{1+o(1)}$, construct at most
$M^{1/3+o(1)}$ children of length $M^{2/3+o(1)}$.  Because the children
cannot explicitly retain all source variables, this lemma must prove actual
variable elimination, decomposition, or a semantic quotient.  It cannot be
obtained by deleting duplicate clauses alone.

\paragraph{Shared materialization lemma.}
Given the certified cover, stream all ordinary residual encodings in time
linear up to subpolynomial factors in their aggregate length.  A construction
that rescans an $M$-bit parent separately for each of $q$ branches costs
$qM$, even if $q\max_i|F_i|\le M^{1+o(1)}$.  Source indexing, persistent
clause views, and delta application can solve the writing problem only after
the residuals themselves are known to be short.

The first lemma is an attractive dense research target because it is not
conditionally forbidden by current compression lower bounds.  The second is
the more severe all-input requirement.  Proving the first without the second
would yield a restricted theorem, not the WRAPPER contract used by the proof
spine.

\section{Transfer matrix}

\renewcommand{\arraystretch}{1.23}
\begin{longtable}{@{}p{0.19\textwidth}p{0.22\textwidth}p{0.25\textwidth}p{0.25\textwidth}@{}}
\toprule
Mechanism & Exact theorem supplied & Quantitative cost or hypothesis & Transfer to the target \\
\midrule
\endhead
Dell--van Melkebeek & No $O(n^b)$ queries of length $O(n^c)$ for $b+c<d$, unless $\mathrm{coNP}\subseteq\mathrm{NP/poly}$ & Polynomial-time reduction; $n$-variable $d$-CNF & Places dense $3$-SAT target exactly at $b+c=3$; no construction \\
IPZ/CIP sparsification & OR of $2^{\varepsilon n}$ sparse $d$-CNFs & Child factor $(d/\varepsilon)^{3d}n$ in the quantitative bound & Polynomial fanout forces a child bound far above $n^2$ \\
2026 measure-and-conquer sparsification & Short proof matching $(d/\varepsilon)^{O(d)}$ occurrence dependence & Fixed $d$, arbitrary positive $\varepsilon$ & Improves proof and dependence, not the polynomial-fanout interpolation \\
Backdoor set & Exact $2^k$ restriction cover into a tractable base & Detection $O(2^kN)$ for Horn and $O(3^kN)$ for $2$CNF & No all-input small $k$; residuals need not shrink; detection exceeds linear target \\
Backdoor DNF & Exact tautological cover by heterogeneous partial assignments & FPT in term count for named base classes & Best syntax for the new lemma; existence, residual size, and near-linear acquisition open \\
Recursive backdoor & Exact recursive assignment/component decomposition & Double-exponential parameter dependence for the empty base & Models recursive WRAPPER but does not meet growing-parameter budget \\
NAE-SAT kernel & $O(n^{d-1})$ clauses for fixed $d\ge4$ & Algebraic dependence rule; conditionally optimal & Strong positive clue; does not transfer to ordinary CNF or requested exponent \\
Representative sets & Preserve all matroid-independent future extensions & Size $\binom{p+q}{p}$ and explicit-family cost & Requires a low-rank SAT extension invariant not supplied by the theorem \\
Universal sets and splitters & Cover every pattern on every chosen $k$-set & Polynomial ambient size only for logarithmic $k$ & Useful after a small relevance theorem; general witness dimension remains large \\
Covering codes & Deterministic local-search coverage & Exponential time in radius or $n$ & Exact algorithmic cover, not a polynomial short-instance cover \\
Split-and-list & Exact meet-in-the-middle compatibility representation & $2^{n/2}$ lists in the basic reduction & Shared combination clue; frontier remains exponential \\
Decision diagrams/DNNF & Canonicalization and exact merge after compilation & FPT under width promises; exponential in general & Valid portfolio member; general acquisition fails \\
Semantic hashing & Exact candidate validation inside compiled decision-DNNF & Polynomial in compiled representation & Minimizes after acquisition; does not bound compilation size \\
OR-distillation lower bounds & Generic many-to-one compression causes complexity collapse & Output independent of number of instances & Warns against one-output shortcut; does not exclude aggregate-linear multi-query target \\
P-selectors & Self-reducible and p-selective languages are in $P$ & Qualitative polynomial time & Shows successful-route choice is solver-strength; no fine-grained exponent \\
Williams self-improvement & Subtracts a fixed large-circuit exponential rate & Fixed $\alpha,\beta$ large-circuit hypotheses & Architectural precedent; no ordinary-CNF all-input cover compiler \\
\bottomrule
\end{longtable}

\section{Final status}

The literature search yields a sharper final statement than either ``known
compression proves it'' or ``compression lower bounds refute it.''  Both are
incorrect.  The balanced target is a conservation frontier.  On dense
$3$-SAT it saturates the Dell--van Melkebeek exponent inequality.  On sparse
inputs it requires genuine variable or semantic reduction that clause
sparsification does not provide.

The constructive bridge worth pursuing is Definition~\ref{def:scbdnf}, with
three nonnegotiable additions: a locally checkable coverage certificate,
individual residual size bounds, and aggregate discovery plus writing cost.
Backdoor DNFs, bounded-width compilation, representative families, and the
double-wrapper separation each implement part of this specification.  No
primary source reviewed here proves favorable coverage on every CNF.

Accordingly, the final link is
\[
 \boxed{\begin{minipage}{0.82\textwidth}
 Prove a uniform, source-checkable, size-charged restriction cover for every
 CNF, with $q\le M^{1/3+o(1)}$, every explicit child of length at most
 $M^{2/3+o(1)}$, and total discovery, materialization, replay, and combination
 work $M^{1+o(1)}$; or replace it by a model-matched theorem giving the same
 outer-solver accounting.
 \end{minipage}}
\]
This statement is \Open.  Therefore the present literature does not complete
the proof spine and does not establish $P\ne NP$.

\begin{thebibliography}{99}

\bibitem{dellvanmelkebeek2014}
Holger Dell and Dieter van Melkebeek,
\emph{Satisfiability Allows No Nontrivial Sparsification Unless the
Polynomial-Time Hierarchy Collapses},
Journal of the ACM 61(4), Article 23, 2014; STOC 2010 version.
\url{https://pages.cs.wisc.edu/~dieter/Papers/sparsification-full.pdf}

\bibitem{ipz2001}
Russell Impagliazzo, Ramamohan Paturi, and Francis Zane,
\emph{Which Problems Have Strongly Exponential Complexity?},
Journal of Computer and System Sciences 63(4):512--530, 2001.
\url{https://cseweb.ucsd.edu/~russell/ipz.pdf}

\bibitem{calabroimpagliazzopaturi2010}
Chris Calabro, Russell Impagliazzo, and Ramamohan Paturi,
\emph{On the Exact Complexity of Evaluating Quantified $k$-CNF},
IPEC 2010, Lecture Notes in Computer Science 6478, pages 50--59.
\url{https://cseweb.ucsd.edu/~paturi/myPapers/pubs/CalabroImpagliazzoPaturi_2010_ipec.pdf}

\bibitem{lokshtanov2026}
Daniel Lokshtanov, Saket Saurabh, and Jie Xue,
\emph{The Sparsification Lemma via Measure and Conquer},
SOSA 2026, pages 291--297.
\url{https://doi.org/10.1137/1.9781611978964.22}

\bibitem{santhanamsrinivasan2011}
Rahul Santhanam and Srikanth Srinivasan,
\emph{On the Limits of Sparsification},
ECCC TR11-131, 2011.
\url{https://eccc.weizmann.ac.il/report/2011/131/download/}

\bibitem{ordyniakschidlerszeider2021}
Sebastian Ordyniak, Andre Schidler, and Stefan Szeider,
\emph{Backdoor DNFs},
IJCAI 2021, pages 1403--1409.
\url{https://www.ijcai.org/proceedings/2021/0194.pdf}

\bibitem{nishimuraragdeszeider2004}
Naomi Nishimura, Prabhakar Ragde, and Stefan Szeider,
\emph{Detecting Backdoor Sets with Respect to Horn and Binary Clauses},
SAT 2004.
\url{https://www.satisfiability.org/SAT04/programme/51.pdf}

\bibitem{mahlmannsiebertzvigny2021}
Nikolas Mahlmann, Sebastian Siebertz, and Alexandre Vigny,
\emph{Recursive Backdoors for SAT},
MFCS 2021, Article 73.
\url{https://drops.dagstuhl.de/entities/document/10.4230/LIPIcs.MFCS.2021.73}

\bibitem{fortnowsanthanam2011}
Lance Fortnow and Rahul Santhanam,
\emph{Infeasibility of Instance Compression and Succinct PCPs for NP},
Journal of Computer and System Sciences 77(1):91--106, 2011.
\url{https://eccc.weizmann.ac.il/eccc-reports/2007/TR07-096/Paper.pdf}

\bibitem{drucker2015}
Andrew Drucker,
\emph{New Limits to Classical and Quantum Instance Compression},
SIAM Journal on Computing 44(5):1443--1479, 2015.
\url{https://eccc.weizmann.ac.il/report/2012/112/}

\bibitem{bodlaenderjansenkratsch2014}
Hans L. Bodlaender, Bart M. P. Jansen, and Stefan Kratsch,
\emph{Kernelization Lower Bounds by Cross-Composition},
SIAM Journal on Discrete Mathematics 28(1):277--305, 2014.
\url{https://arxiv.org/abs/1206.5941}

\bibitem{jansenpieterse2017}
Bart M. P. Jansen and Astrid Pieterse,
\emph{Sparsification Upper and Lower Bounds for Graph Problems and
Not-All-Equal SAT},
Algorithmica 79:3--28, 2017.
\url{https://arxiv.org/abs/1509.07437}

\bibitem{fominrepresentative2016}
Fedor V. Fomin, Daniel Lokshtanov, Fahad Panolan, and Saket Saurabh,
\emph{Efficient Computation of Representative Families with Applications in
Parameterized and Exact Algorithms},
Journal of the ACM 63(4), Article 29, 2016.
\url{https://arxiv.org/abs/1304.4626}

\bibitem{naorschulmansrinivasan1995}
Moni Naor, Leonard J. Schulman, and Aravind Srinivasan,
\emph{Splitters and Near-Optimal Derandomization},
FOCS 1995, pages 182--191.
\url{https://www.wisdom.weizmann.ac.il/~naor/PAPERS/splitters.pdf}

\bibitem{moserscheder2011}
Robin A. Moser and Dominik Scheder,
\emph{A Full Derandomization of Sch\"oning's $k$-SAT Algorithm},
STOC 2011, pages 245--252.
\url{https://arxiv.org/abs/1008.4067}

\bibitem{patrascuwilliams2010}
Mihai P\u{a}tra\c{s}cu and Ryan Williams,
\emph{On the Possibility of Faster SAT Algorithms},
SODA 2010, pages 1065--1075.
\url{https://people.csail.mit.edu/rrw/sat-talg.pdf}

\bibitem{bryant1986}
Randal E. Bryant,
\emph{Graph-Based Algorithms for Boolean Function Manipulation},
IEEE Transactions on Computers C-35(8):677--691, 1986.
\url{https://www.cs.cmu.edu/~bryant/pubdir/ieeetc86.pdf}

\bibitem{bierelagniezlonca2026}
Armin Biere, Jean-Marie Lagniez, and Emmanuel Lonca,
\emph{Efficient Minimization of Decision-DNNF Circuits via Semantic Hashing
and Provenance Tracking},
IJCAI 2026.
\url{https://www.ijcai.org/proceedings/2026/420}

\bibitem{capelli2026}
Florent Capelli, YooJung Choi, Stefan Mengel, Mart\'in Mu\~noz, and
Guy Van den Broeck,
\emph{A Canonical Generalization of OBDD},
SAT 2026, Article 10.
\url{https://drops.dagstuhl.de/entities/document/10.4230/LIPIcs.SAT.2026.10}

\bibitem{bovaetal2016}
Simone Bova, Florent Capelli, Stefan Mengel, and Friedrich Slivovsky,
\emph{A Strongly Exponential Separation of DNNFs from CNF Formulas},
ISAAC 2016.
\url{https://arxiv.org/abs/1411.1995}

\bibitem{amarillimonetsenellart2018}
Antoine Amarilli, Mika\"el Monet, and Pierre Senellart,
\emph{Connecting Width and Structure in Knowledge Compilation},
ICDT 2018, Article 6.
\url{https://drops.dagstuhl.de/entities/document/10.4230/LIPIcs.ICDT.2018.6}

\bibitem{valiantvazirani1986}
Leslie G. Valiant and Vijay V. Vazirani,
\emph{NP Is as Easy as Detecting Unique Solutions},
Theoretical Computer Science 47:85--93, 1986.
\url{https://doi.org/10.1016/0304-3975(86)90135-0}

\bibitem{buhrmantorenvliet1996}
Harry Buhrman and Leen Torenvliet,
\emph{P-Selective Self-Reducible Sets: A New Characterization of P},
Journal of Computer and System Sciences 53(2):210--217, 1996.
\url{https://doi.org/10.1006/jcss.1996.0062}

\bibitem{williams2024}
R. Ryan Williams,
\emph{Self-Improvement for Circuit-Analysis Problems},
STOC 2024, pages 1374--1385.
\url{https://people.csail.mit.edu/rrw/stoc-self-improvement.pdf}

\end{thebibliography}

\end{document}
% END IMMUTABLE RESEARCH REPORT::FINAL-LINK-LIT::2026-10-08_003222
